\chapter{Cadre général du projet}
\label{chap:cadre}


Ce premier chapitre situe le projet dans son double contexte : celui de l'organisme
qui l'a accueilli d'une part, celui de la médecine du sommeil d'autre part. Nous
présentons d'abord le Centre National de l'Informatique et la place qu'y occupe ce
travail. Nous exposons ensuite les notions médicales indispensables à la
compréhension du mémoire : l'architecture du sommeil, l'examen polysomnographique et
le syndrome d'apnées obstructives du sommeil. De ce contexte émerge une problématique
en trois volets, que nous confrontons aux solutions existantes, commerciales comme
académiques. Nous en déduisons enfin le périmètre de la solution proposée et la
méthodologie retenue pour la construire.

% =============================================================================
\section{Présentation de l'organisme d'accueil}
% =============================================================================

\subsection{Le Centre National de l'Informatique}

Le Centre National de l'Informatique (CNI) est un établissement public à caractère
non administratif, doté de la personnalité civile et de l'autonomie financière. Il a
été créé par la loi n\textsuperscript{o} 75-83 du 30 décembre 1975, dont les articles
35 à 42 fixent son statut et ses attributions. Placé sous la tutelle du ministère
chargé des technologies de la communication, il est certifié ISO 9001 version 2015.

Le CNI constitue le principal appui technique des structures publiques tunisiennes
pour la conception, le déploiement et l'exploitation de leurs systèmes d'information.
À ce titre, il intervient aussi bien en maîtrise d'ouvrage déléguée qu'en réalisation
directe, et opère une infrastructure d'hébergement mutualisée au service de
l'administration.

\subsection{Missions}

Les missions du CNI couvrent l'ensemble du cycle de vie d'un système d'information
public :

\begin{itemize}[itemsep=2pt,topsep=4pt]
  \item \textbf{Maîtrise d'ouvrage déléguée} --- pilotage de projets pour le compte
        des administrations bénéficiaires.
  \item \textbf{Études et conseil} --- audits, études d'opportunité, schémas
        directeurs, rédaction de cahiers des charges, conseil en réseaux et en
        sécurité des systèmes d'information.
  \item \textbf{Développement} --- conception, réalisation et maintenance
        applicative, accompagnée de la formation des utilisateurs aux systèmes
        livrés.
  \item \textbf{Hébergement} --- exploitation de serveurs, d'applications et de
        données pour le compte des organismes publics.
  \item \textbf{Continuité d'activité} --- solutions de sauvegarde, plans de reprise
        après sinistre et plans de continuité.
  \item \textbf{Formation et certification} --- formations continues et certifiantes,
        séminaires techniques.
  \item \textbf{Déploiement et assistance} --- mise en production des applications,
        support sur site ou à distance.
\end{itemize}

\subsection{Repères historiques}

Le CNI a accompagné toutes les vagues d'informatisation de l'administration
tunisienne, depuis les traitements centralisés des années 1970 jusqu'aux services
numériques contemporains. La figure~\ref{fig:cni-timeline} en retrace les principaux
jalons.

\begin{figure}[H]
\centering
\begin{tikzpicture}[x=1cm,y=1cm,
    jalon/.style   = {align=center,font=\scriptsize,text width=2.6cm},
    millesime/.style = {font=\footnotesize\bfseries,text=hypnoblue},
  ]

  % Axe temporel
  \draw[line width=1.1pt,hypnoblue,-{Latex[length=3mm]}] (-0.4,0) -- (15,0);

  % Jalons portés au-dessus de l'axe
  \foreach \px/\yr/\txt in {
      0/1975/{Création du CNI\\par la loi 75-83},
      6/1992/{Élargissement de\\l'offre de formation},
      12/2007/{Réseau national,\\e-gouvernement, backup}}
  {
    \draw[line width=1pt,hypnoblue] (\px,-0.18) -- (\px,0.18);
    \fill[hypnored] (\px,0) circle (2.4pt);
    \node[jalon,anchor=south]     at (\px, 0.62) {\txt};
    \node[millesime,anchor=north] at (\px,-0.24) {\yr};
  }

  % Jalons portés sous l'axe
  \foreach \px/\yr/\txt in {
      3/1978/{Prise en charge de la\\formation en informatique},
      9/1997/{Reclassement en\\établissement public},
      14.4/{Auj.}/{Pôle d'excellence\\des TIC publiques}}
  {
    \draw[line width=1pt,hypnoblue] (\px,-0.18) -- (\px,0.18);
    \fill[hypnored] (\px,0) circle (2.4pt);
    \node[jalon,anchor=north]     at (\px,-0.62) {\txt};
    \node[millesime,anchor=south] at (\px, 0.24) {\yr};
  }

\end{tikzpicture}
\caption[Chronologie du Centre National de l'Informatique]%
        {Principaux jalons de l'évolution du Centre National de l'Informatique.}
\label{fig:cni-timeline}
\end{figure}

\subsection{Positionnement du projet}

Le projet \hypnoria{} s'inscrit dans les axes « Études et conseil » et
« Développement » du CNI. Il relève d'une démarche exploratoire : appliquer les
techniques d'apprentissage automatique récentes à un domaine — l'aide au diagnostic
médical — où l'administration tunisienne dispose de gisements de données mais de peu
d'outils d'exploitation. Le stage a bénéficié d'un environnement de développement
complet, d'un accès à des ressources de calcul accélérées et d'un encadrement
technique régulier.

Trois caractéristiques du contexte CNI ont orienté les choix techniques du projet :

\begin{itemize}[itemsep=2pt,topsep=4pt]
  \item \textbf{La souveraineté des données.} Les données de santé sont sensibles.
        L'architecture retenue sépare strictement les métadonnées cliniques, stockées
        dans une base relationnelle maîtrisée, des fichiers volumineux, confiés à un
        stockage objet dont l'accès n'est possible que par des liens signés à durée
        limitée.
  \item \textbf{La reproductibilité du déploiement.} Le CNI exploite des systèmes sur
        une longue durée. L'ensemble de la plateforme est donc conteneurisé et
        démarrable en une seule commande, sans dépendance à une configuration de
        poste particulière.
  \item \textbf{La traçabilité.} Toute action sensible est journalisée avec son
        auteur, son horodatage et son adresse d'origine, et le journal est exportable
        en vue d'un contrôle de conformité.
\end{itemize}

% =============================================================================
\section{Contexte médical}
% =============================================================================

\subsection{Le sommeil, un processus cyclique et structuré}

Le sommeil n'est pas un état homogène de repos mais une succession organisée d'états
physiologiques distincts, qui se répètent au cours de la nuit selon des cycles d'une
durée moyenne de 90 minutes. On distingue deux grandes familles d'états : le sommeil
sans mouvements oculaires rapides, ou \emph{NREM}, subdivisé en stades de profondeur
croissante, et le sommeil paradoxal, ou \emph{REM}, caractérisé par une activité
cérébrale proche de l'éveil associée à une atonie musculaire complète.

Cette organisation n'est pas anecdotique : elle porte l'essentiel de la valeur
diagnostique de l'examen. La proportion de sommeil profond conditionne la
récupération physique, celle de sommeil paradoxal la consolidation mnésique, et la
fragmentation de l'ensemble constitue le marqueur principal de nombreuses pathologies
respiratoires nocturnes.

\subsection{Les stades du sommeil selon l'AASM}

La classification de référence est celle de l'\emph{American Academy of Sleep
Medicine}~\cite{aasm2020}, qui a succédé en 2007 aux règles de Rechtschaffen et
Kales~\cite{rechtschaffen1968}. Elle définit cinq états, décrits dans le
tableau~\ref{tab:stades}, et impose que le scorage se fasse par fenêtres, ou
\emph{époques}, de 30 secondes : à chaque époque est attribué un et un seul stade,
celui qui occupe la majorité de la fenêtre.

La différence majeure entre les deux référentiels tient à la fusion des anciens
stades 3 et 4 en un unique stade N3. Cette convention sera reprise dans tout ce
mémoire, y compris lors du traitement des annotations de référence des cohortes
utilisées, qui suivent encore la nomenclature antérieure.

\begin{table}[H]
\centering
\caption{Les cinq états de vigilance définis par le référentiel AASM}
\label{tab:stades}
\small
\begin{tabular}{@{}lL{4.6cm}L{5.4cm}c@{}}
\toprule
\textbf{Stade} & \textbf{Signature électrophysiologique} & \textbf{Éléments
caractéristiques} & \textbf{Part typique} \\
\midrule
W (éveil) & Rythme alpha postérieur, activité bêta, tonus musculaire élevé &
Mouvements oculaires volontaires et clignements & --- \\
\addlinespace[2pt]
N1 & Ralentissement du rythme de fond, ondes thêta & Mouvements oculaires lents,
transition instable & \SI{5}{\percent} \\
\addlinespace[2pt]
N2 & Activité thêta dominante & Fuseaux du sommeil et complexes K & \SI{45}{\percent}
à \SI{55}{\percent} \\
\addlinespace[2pt]
N3 & Ondes lentes de grande amplitude, sur plus de \SI{20}{\percent} de l'époque &
Sommeil profond, seuil d'éveil élevé & \SI{15}{\percent} à \SI{25}{\percent} \\
\addlinespace[2pt]
REM & Rythme rapide et désynchronisé, proche de l'éveil & Mouvements oculaires
rapides, atonie musculaire & \SI{20}{\percent} à \SI{25}{\percent} \\
\bottomrule
\end{tabular}
\end{table}

Deux points méritent d'être soulignés, car ils conditionnent la difficulté de la
tâche d'automatisation traitée dans ce mémoire.

D'une part, la distribution des stades est fortement déséquilibrée. Le N2 représente
à lui seul environ la moitié d'une nuit, tandis que le N1 en occupe à peine un
vingtième. Tout modèle entraîné sans précaution sur ces données apprendra
préférentiellement la classe majoritaire.

D'autre part, le stade N1 est un état de transition dont les frontières sont
intrinsèquement floues. Il ne possède aucun marqueur graphoélément propre : il se
définit par la disparition du rythme alpha et l'absence des fuseaux du N2, c'est-à-dire
négativement. C'est la classe sur laquelle les experts humains eux-mêmes s'accordent
le moins.

\subsection{L'hypnogramme}

La succession des stades attribués aux époques successives constitue
l'\emph{hypnogramme}, représentation graphique synthétique d'une nuit
(figure~\ref{fig:hypnogramme}). C'est le document de référence du médecin du sommeil :
il condense environ mille décisions élémentaires en une seule courbe, à partir de
laquelle se calculent toutes les métriques cliniques.

\begin{figure}[H]
\centering
\begin{tikzpicture}[x=1.02cm,y=1cm,
    annot/.style = {font=\scriptsize,align=center},
    fleche/.style= {-{Latex[length=2mm]},line width=0.55pt},
  ]

  % ── Bandes signalant les épisodes de sommeil paradoxal ────────────────────
  \begin{scope}[on background layer]
    \foreach \a/\b in {2.25/2.625, 4.5/5.025, 7.125/7.875, 9.0/9.9, 10.95/12.0}
      \fill[stREM!14] (\a,-3.95) rectangle (\b,0.30);
  \end{scope}

  % ── Lignes de niveau ──────────────────────────────────────────────────────
  \foreach \y in {0,-0.9,-1.8,-2.7,-3.6}
    \draw[hypnogray!22,line width=0.4pt] (0,\y) -- (12.3,\y);

  % ── Étiquettes des stades ─────────────────────────────────────────────────
  \node[left,font=\scriptsize\bfseries,text=stWake] at (-0.12,0)    {W};
  \node[left,font=\scriptsize\bfseries,text=stREM]  at (-0.12,-0.9) {REM};
  \node[left,font=\scriptsize\bfseries,text=stN1]   at (-0.12,-1.8) {N1};
  \node[left,font=\scriptsize\bfseries,text=stN2]   at (-0.12,-2.7) {N2};
  \node[left,font=\scriptsize\bfseries,text=stN3]   at (-0.12,-3.6) {N3};

  % ── Tracé ─────────────────────────────────────────────────────────────────
  \draw[line width=1.15pt,hypnoblue]
    (0,0) -- (0.45,0) -- (0.45,-1.8) -- (0.675,-1.8) -- (0.675,-2.7) --
    (1.275,-2.7) -- (1.275,-3.6) -- (1.95,-3.6) -- (1.95,-2.7) --
    (2.25,-2.7) -- (2.25,-0.9) -- (2.625,-0.9) -- (2.625,-1.8) --
    (2.775,-1.8) -- (2.775,-2.7) -- (3.45,-2.7) -- (3.45,-3.6) --
    (4.2,-3.6) -- (4.2,-2.7) -- (4.5,-2.7) -- (4.5,-0.9) --
    (5.025,-0.9) -- (5.025,0) -- (5.175,0) -- (5.175,-1.8) --
    (5.325,-1.8) -- (5.325,-2.7) -- (6.3,-2.7) -- (6.3,-3.6) --
    (6.75,-3.6) -- (6.75,-2.7) -- (7.125,-2.7) -- (7.125,-0.9) --
    (7.875,-0.9) -- (7.875,0) -- (8.025,0) -- (8.025,-1.8) --
    (8.175,-1.8) -- (8.175,-2.7) -- (9.0,-2.7) -- (9.0,-0.9) --
    (9.9,-0.9) -- (9.9,0) -- (10.125,0) -- (10.125,-1.8) --
    (10.35,-1.8) -- (10.35,-2.7) -- (10.95,-2.7) -- (10.95,-0.9) --
    (12.0,-0.9) -- (12.0,0) -- (12.3,0);

  % ── Axe des temps ─────────────────────────────────────────────────────────
  \draw[hypnogray,line width=0.6pt] (0,-4.05) -- (12.3,-4.05);
  \foreach \h in {0,1,2,3,4,5,6,7,8} {
    \pgfmathsetmacro{\xx}{\h*1.5}
    \draw[hypnogray,line width=0.5pt] (\xx,-4.05) -- (\xx,-4.22);
    \node[below,font=\scriptsize,text=hypnogray,inner sep=1.5pt] at (\xx,-4.22) {\h\,h};
  }

  % ── Annotation : sommeil profond ──────────────────────────────────────────
  \node[annot,text=hypnored,text width=3.9cm,anchor=north] at (2.6,-4.85)
       {sommeil profond concentré\\dans le premier tiers de la nuit};
  \draw[fleche,hypnored] (1.75,-4.80) -- (1.62,-3.72);

  % ── Annotation : sommeil paradoxal ────────────────────────────────────────
  \node[annot,text=stREM,text width=4.3cm,anchor=south] at (9.7,1.05)
       {épisodes de sommeil paradoxal\\de plus en plus longs};
  \draw[fleche,stREM] (10.75,1.00) -- (11.45,-0.78);

\end{tikzpicture}
\vspace{0.35cm}
\caption[Hypnogramme d'une nuit physiologique]%
        {Hypnogramme schématique d'une nuit physiologique de huit heures. Les bandes
         vertes signalent les épisodes de sommeil paradoxal.}
\label{fig:hypnogramme}
\end{figure}

Deux régularités physiologiques sont directement lisibles sur cette figure, et le
modèle devra les retrouver : le sommeil profond se concentre dans le premier tiers de
la nuit, tandis que les épisodes de sommeil paradoxal s'allongent progressivement
jusqu'au réveil.

\subsection{La polysomnographie}

L'hypnogramme n'est pas mesuré, il est \emph{interprété}. Il découle d'un examen
appelé polysomnographie (PSG), qui enregistre simultanément plusieurs signaux
physiologiques pendant une nuit complète. Le tableau~\ref{tab:derivations} décrit les
dérivations minimales requises par le référentiel AASM pour établir un scorage, qui
sont précisément celles exploitées dans ce travail.

\begin{table}[H]
\centering
\caption{Dérivations minimales d'un montage polysomnographique de scorage}
\label{tab:derivations}
\small
\begin{tabular}{@{}llL{7.2cm}@{}}
\toprule
\textbf{Signal} & \textbf{Nombre} & \textbf{Apport au scorage} \\
\midrule
Électroencéphalogramme (EEG) & 2 & Rythme de fond, fuseaux, complexes K et ondes
lentes : c'est le signal qui porte la décision \\
\addlinespace[2pt]
Électro-oculogramme (EOG) & 2 & Distingue les mouvements oculaires lents de
l'endormissement des mouvements rapides du sommeil paradoxal \\
\addlinespace[2pt]
Électromyogramme (EMG) mentonnier & 1 & Le tonus musculaire sépare le sommeil
paradoxal du sommeil lent : son effondrement signe le REM \\
\bottomrule
\end{tabular}
\end{table}

Un montage complet comporte en pratique d'autres capteurs — flux nasal, sangles
thoraciques et abdominales, oxymétrie de pouls, électrocardiogramme, capteur de
position — destinés à la caractérisation des événements respiratoires. Le scorage des
stades, lui, ne requiert que les cinq dérivations ci-dessus.

Ce point est structurant pour notre travail : il rend possible un système de
classification qui n'exige que cinq canaux, et donc compatible avec des montages
allégés. Nous évaluerons d'ailleurs une configuration réduite à deux dérivations EEG,
correspondant aux dispositifs d'enregistrement ambulatoire.

\subsection{Le syndrome d'apnées obstructives du sommeil}

Le syndrome d'apnées obstructives du sommeil (SAOS) résulte d'obstructions répétées
des voies aériennes supérieures pendant le sommeil. Chaque obstruction provoque une
désaturation en oxygène et un micro-éveil, invisible pour le patient mais destructeur
pour l'architecture du sommeil.

Sa prévalence est considérable : une analyse de la littérature publiée en 2019 estime
à environ 936 millions le nombre d'adultes de 30 à 69 ans présentant un SAOS léger à
sévère dans le monde, dont près de 425 millions sous une forme modérée à
sévère~\cite{benjafield2019}. Non traité, le syndrome est associé à une
surmorbidité cardiovasculaire documentée~\cite{marin2005}. Il reste pourtant
largement sous-diagnostiqué~\cite{young1997}, principalement en raison du coût et de
la disponibilité limitée de l'examen de référence.

La sévérité s'exprime par l'index d'apnées-hypopnées (IAH), nombre moyen d'événements
respiratoires par heure de sommeil. Le référentiel AASM~\cite{kapur2017} en définit
quatre classes, reprises au tableau~\ref{tab:iah} et qui constituent les classes cibles
du second modèle développé dans ce mémoire.

\begin{table}[H]
\centering
\caption{Classes de sévérité du SAOS selon l'index d'apnées-hypopnées}
\label{tab:iah}
\small
\begin{tabular}{@{}lcL{7.4cm}@{}}
\toprule
\textbf{Classe} & \textbf{IAH (évén./h)} & \textbf{Conduite habituelle} \\
\midrule
Normal   & $< 5$          & Pas de traitement spécifique \\
Léger    & $5$ à $< 15$   & Mesures hygiéno-diététiques, surveillance \\
Modéré   & $15$ à $< 30$  & Traitement discuté selon la symptomatologie \\
Sévère   & $\geq 30$      & Traitement par pression positive continue indiqué \\
\bottomrule
\end{tabular}
\end{table}

Le point essentiel pour la suite est le suivant : le SAOS ne se contente pas de
provoquer des événements respiratoires, il \emph{déforme l'architecture du sommeil}.
Les micro-éveils répétés fragmentent la nuit, suppriment le sommeil profond et
raccourcissent les épisodes paradoxaux. Cette empreinte est mesurable sur le seul
hypnogramme, indépendamment des capteurs respiratoires — c'est l'hypothèse sur
laquelle repose le second volet de ce projet.

% =============================================================================
\section{Problématique}
% =============================================================================

\subsection{Le coût du scorage manuel}

Une nuit d'enregistrement de huit heures produit environ un millier d'époques de
30 secondes. Chacune doit être examinée, comparée à ses voisines et classée par un
technicien du sommeil qualifié. L'opération représente plusieurs heures de travail
expert par patient, auxquelles s'ajoutent l'analyse des événements respiratoires et
la rédaction du compte rendu.

Ce coût se traduit directement en délais. Dans un service dont la capacité de lecture
est saturée, le facteur limitant n'est plus le nombre de lits d'enregistrement mais le
temps d'expertise disponible pour interpréter ce qu'ils produisent. C'est un des
mécanismes du sous-diagnostic évoqué plus haut.

\subsection{La variabilité inter-scoreur}

Le scorage manuel n'est pas seulement coûteux, il est également imparfaitement
reproductible. Le programme de fiabilité inter-scoreurs de l'AASM, qui compare les
annotations de plusieurs milliers de scoreurs sur des enregistrements communs, mesure
un accord global de l'ordre de \SI{83}{\percent}~\cite{rosenberg2013}. D'autres
travaux multicentriques rapportent des ordres de grandeur comparables, avec des écarts
marqués selon le stade~\cite{danker2009}.

Cette dispersion n'est pas uniforme. L'accord est excellent sur l'éveil et le sommeil
paradoxal, dont les signatures sont franches, et nettement plus faible sur le N1,
dont nous avons vu qu'il se définit par défaut. Deux conséquences en découlent.

La première est méthodologique : les annotations de référence utilisées pour entraîner
et évaluer un modèle ne sont pas une vérité absolue mais l'opinion d'un scoreur. Un
modèle qui atteindrait un accord parfait avec un scoreur donné serait en désaccord
avec un autre. L'accord inter-expert constitue donc un plafond de performance
réaliste, et c'est à cette aune qu'il faut lire les résultats du
chapitre~\ref{chap:realisation-ia}.

La seconde est clinique : elle légitime l'automatisation. Un modèle, quelles que
soient ses limites, présente au moins la propriété d'être parfaitement reproductible.
Appliqué deux fois au même enregistrement, il rend deux fois la même réponse.

\subsection{L'opacité des modèles d'apprentissage profond}

Depuis 2017, les réseaux de neurones profonds atteignent sur les bases publiques des
performances comparables à celles d'un scoreur humain~\cite{supratak2017,perslev2021}.
Le verrou n'est donc plus principalement la précision : il est l'adoption.

Un modèle d'apprentissage profond fonctionne comme une boîte noire. Il produit une
classe sans énoncer ce qui l'a déterminée. Or le praticien engage sa responsabilité
sur le diagnostic qu'il signe, et ne peut ni valider ni contester une décision dont
il ignore les motifs. À ce frein pratique s'ajoute un cadre réglementaire européen
qui reconnaît un droit à l'explication pour les décisions automatisées ayant un effet
significatif sur les personnes~\cite{goodman2017}.

L'explicabilité n'est donc pas un ornement de la prédiction : elle en est la
condition d'usage clinique.

\subsection{Synthèse}

Ces trois constats convergent vers une formulation unique. La médecine du sommeil a
besoin d'un système qui soit à la fois \textbf{rapide}, pour absorber le volume
d'examens ; \textbf{reproductible}, pour lever la variabilité inter-scoreur ; et
\textbf{transparent}, pour que le praticien puisse s'en porter garant. C'est cette
triple exigence que \hypnoria{} cherche à satisfaire.

% =============================================================================
\section{Étude de l'existant}
% =============================================================================

\subsection{Les solutions commerciales}

Le marché de l'analyse polysomnographique est dominé par des logiciels propriétaires
adossés à du matériel d'acquisition : \emph{Noxturnal} de Nox Medical, \emph{Sleepware
G3} de Philips Respironics, \emph{ProFusion PSG} de Compumedics ou encore
\emph{SleepWorks} de Natus. Une seconde génération d'acteurs, dont \emph{EnsoSleep}
d'EnsoData, propose un scorage automatique fondé sur l'apprentissage profond.

Ces outils sont matures, validés cliniquement et parfaitement adaptés à leur usage.
Ils présentent néanmoins, du point de vue de notre problématique, quatre limites
communes.

\begin{itemize}[itemsep=2pt,topsep=4pt]
  \item \textbf{L'adhérence matérielle.} La plupart sont conçus pour exploiter les
        enregistrements produits par le matériel du même fabricant, ce qui contraint
        les laboratoires disposant d'un parc hétérogène.
  \item \textbf{Le coût de licence.} Le modèle économique, calibré pour des
        structures hospitalières occidentales, reste hors de portée de nombreux
        centres.
  \item \textbf{L'installation locale.} L'architecture est le plus souvent celle d'un
        logiciel de poste, peu compatible avec le travail à distance et la
        collaboration entre praticiens de sites différents.
  \item \textbf{L'absence d'explication.} Lorsqu'un scorage automatique est proposé,
        il l'est sous forme d'un résultat à valider, sans exposition des facteurs qui
        l'ont produit.
\end{itemize}

\subsection{Les travaux académiques}

La littérature scientifique a produit une succession d'architectures de plus en plus
performantes. \emph{DeepSleepNet}~\cite{supratak2017} combine une extraction
convolutive multi-échelle et une modélisation séquentielle récurrente sur un unique
canal EEG. \emph{TinySleepNet}~\cite{supratak2020} en propose une version allégée à
destination des dispositifs embarqués. \emph{U-Sleep}~\cite{perslev2021} démontre une
robustesse remarquable au montage et à la cohorte. \emph{SleepTransformer}~\cite{phan2022}
introduit un mécanisme d'auto-attention et, fait notable, une quantification de
l'incertitude ainsi qu'une forme d'interprétabilité par les poids d'attention.
\emph{XSleepNet}~\cite{phan2021} exploite conjointement la représentation temporelle
brute et la représentation temps-fréquence.

Ces travaux partagent toutefois une orientation commune : ils visent la performance de
classification, évaluée sur des jeux publics, et s'arrêtent au modèle. Aucun ne livre
la chaîne applicative qui permettrait à un praticien d'en faire usage — gestion des
patients, archivage des examens, production d'un compte rendu, collaboration.

\subsection{Analyse critique et positionnement}

Le tableau~\ref{tab:existant} confronte ces deux familles de solutions aux besoins
identifiés dans la problématique.

\begin{table}[H]
\centering
\caption{Positionnement de \hypnoria{} face aux solutions existantes}
\label{tab:existant}
\small
\begin{tabular}{@{}L{4.5cm}ccc@{}}
\toprule
\textbf{Critère} & \textbf{Suites commerciales} & \textbf{Travaux académiques} &
\textbf{\hypnoria{}} \\
\midrule
Scorage automatique            & partiel  & oui  & oui \\
Indépendance au matériel       & faible   & --   & oui \\
Explicabilité des décisions    & non      & rare & oui \\
Prédiction de la sévérité SAOS & indirecte& non  & oui \\
Application web multi-utilisateurs & non  & non  & oui \\
Dossier patient et historique  & oui      & non  & oui \\
Collaboration entre praticiens & non      & non  & oui \\
Traçabilité et audit           & partielle& non  & oui \\
Coût d'acquisition             & élevé    & nul  & nul \\
\bottomrule
\end{tabular}
\end{table}

Il en ressort un espace non couvert, à l'intersection des deux familles : une
plateforme applicative complète, indépendante du matériel d'acquisition, qui intègre
des modèles de l'état de l'art \emph{et} les explique. C'est cet espace que le projet
occupe.

% =============================================================================
\section{Solution proposée}
% =============================================================================

\subsection{Objectifs}

\hypnoria{} est une plateforme web d'aide à l'interprétation polysomnographique,
destinée aux médecins du sommeil. Elle poursuit quatre objectifs.

\begin{enumerate}[itemsep=3pt,topsep=4pt]
  \item \textbf{Automatiser le scorage des stades} à partir d'un fichier
        d'enregistrement brut au format EDF, avec une exactitude comparable à
        l'accord inter-expert.
  \item \textbf{Prédire la sévérité du SAOS} à partir de l'architecture du sommeil
        et de marqueurs cliniques, sans exiger l'annotation manuelle des événements
        respiratoires.
  \item \textbf{Expliquer chaque prédiction} en exposant les facteurs qui l'ont
        déterminée, afin que le praticien puisse la valider ou la contester.
  \item \textbf{Fournir la chaîne applicative complète} : dossier patient, archivage
        sécurisé, annotation, compte rendu exportable, collaboration et
        administration multi-établissements.
\end{enumerate}

\subsection{Périmètre fonctionnel}

Le système s'articule autour de deux étages d'apprentissage successifs et d'une
plateforme qui les met à disposition.

Le premier étage assure la \textbf{classification des stades}. Trois familles
d'architectures — convolutive, récurrente et attentionnelle — sont entraînées
séparément, puis combinées par un méta-apprenant selon le principe de l'empilement de
modèles~\cite{wolpert1992}. Deux configurations de montage (deux ou cinq dérivations)
et deux granularités de sortie (trois ou cinq stades) sont proposées, afin de couvrir
aussi bien l'enregistrement de laboratoire que le dispositif ambulatoire.

Le second étage assure la \textbf{prédiction de la sévérité}. Il extrait de
l'hypnogramme produit une vingtaine de marqueurs d'architecture du sommeil, les
complète de données cliniques et oxymétriques, et prédit la classe de sévérité. Chaque
prédiction est accompagnée de sa décomposition en contributions par variable, obtenue
par la méthode des valeurs de Shapley~\cite{lundberg2017,lundberg2020}.

La \textbf{plateforme}, enfin, expose ces traitements dans une application web
sécurisée : authentification, séparation des rôles, gestion du cycle de vie des
licences cliniques, registre des patients, archivage des fichiers volumineux sur un
stockage objet, annotation des hypnogrammes, génération de comptes rendus,
discussions entre praticiens et journal d'audit exportable.

\subsection{Contributions}

Les contributions revendiquées par ce travail sont les suivantes :

\begin{itemize}[itemsep=2pt,topsep=4pt]
  \item l'entraînement et l'évaluation systématiques de douze modèles de
        classification couvrant trois architectures, deux configurations de montage
        et deux granularités de sortie, sur une cohorte clinique ;
  \item la construction de quatre ensembles par empilement et l'analyse de leur
        apport réel par rapport aux modèles isolés et au vote moyen ;
  \item une chaîne de prédiction de la sévérité du SAOS fondée sur l'architecture du
        sommeil, explicitement protégée contre les fuites de la variable cible, et
        assortie d'explications individuelles ;
  \item l'intégration de l'ensemble dans une plateforme clinique opérationnelle,
        conteneurisée et traçable ;
  \item une discussion argumentée des limites méthodologiques de l'évaluation, qui
        constitue à nos yeux une contribution à part entière.
\end{itemize}

% =============================================================================
\section{Méthodologie de gestion de projet}
% =============================================================================

\subsection{Le choix d'une démarche agile}

Le projet présentait deux sources d'incertitude difficilement compatibles avec une
planification linéaire. La première tenait aux données : l'accès aux cohortes
cliniques est soumis à une autorisation dont le délai d'instruction n'était pas connu
à l'avance. La seconde tenait aux modèles : rien ne garantissait \emph{a priori}
qu'une architecture donnée atteindrait le niveau de performance visé, ni que
l'empilement apporterait un gain.

Une démarche séquentielle aurait figé des décisions avant de disposer des éléments
permettant de les prendre. Nous avons donc retenu le cadre
Scrum~\cite{schwaber2020}, dont le principe est d'avancer par incréments courts et
d'ajuster la suite à ce que chaque incrément a révélé.

\subsection{Adaptation à un projet individuel}

Scrum suppose une équipe pluridisciplinaire ; le projet est mené par une seule
personne. Le cadre a donc été adapté sans être dénaturé, selon la
figure~\ref{fig:scrum}.

Les trois rôles ont été redistribués : la fonction de \emph{Product Owner} — arbitrage
des priorités et validation de la valeur livrée — a été portée par l'encadrement
professionnel du CNI ; la fonction de \emph{Scrum Master} — respect du cadre et levée
des obstacles — par l'encadrement académique lors des points de suivi ; la fonction
d'équipe de développement par l'étudiant.

La durée d'itération a été portée à quatre semaines plutôt que deux. Ce choix est
directement lié à la nature du travail : l'entraînement d'un modèle profond sur près
de deux cent mille exemples se compte en heures de calcul, et une itération de deux
semaines n'aurait pas permis de produire un incrément évaluable pendant les sprints
de modélisation.

\begin{figure}[H]
\centering
\begin{tikzpicture}[
    node distance=0pt,
    box/.style   = {rectangle,rounded corners=2pt,draw=hypnoblue!55,fill=hypnolight,
                    align=center,font=\scriptsize,minimum height=1.05cm,text width=2.15cm},
    dark/.style  = {rectangle,rounded corners=2pt,draw=hypnored!70,fill=hypnored!8,
                    align=center,font=\scriptsize,minimum height=1.05cm,text width=2.15cm},
    fl/.style    = {-{Latex[length=2.2mm]},hypnoblue!70,line width=0.8pt},
  ]

  \node[box] (backlog) at (0,0)      {\textbf{Product\\ Backlog}\\[1pt] 60 user stories};
  \node[box] (planif)  at (3.3,0)    {\textbf{Planification}\\[1pt] sélection et
                                      estimation};
  \node[dark] (sprint) at (6.6,0)    {\textbf{Sprint}\\[1pt] 4 semaines};
  \node[box] (revue)   at (9.9,0)    {\textbf{Revue}\\[1pt] démonstration à
                                      l'encadrement};
  \node[box] (retro)   at (13.0,0)   {\textbf{Rétrospective}\\[1pt] ajustement du
                                      sprint suivant};

  \draw[fl] (backlog) -- (planif);
  \draw[fl] (planif)  -- (sprint);
  \draw[fl] (sprint)  -- (revue);
  \draw[fl] (revue)   -- (retro);

  % Boucle de retour
  \draw[fl] (retro.south) -- ++(0,-0.85) -- node[below,font=\scriptsize\itshape,
        pos=0.5,hypnogray] {réalimentation du backlog à chaque itération}
        ++(-13.0,0) -- (backlog.south);

  % Incrément
  \node[font=\scriptsize\itshape,hypnogray,align=center] at (6.6,1.15)
       {incrément fonctionnel livré};
  \draw[fl] (sprint.north) -- (6.6,0.95);

\end{tikzpicture}
\vspace{0.6cm}
\caption[Cycle Scrum adapté au projet]%
        {Cycle Scrum adapté à un projet individuel de six mois.}
\label{fig:scrum}
\end{figure}

\subsection{Artefacts, cérémonies et outillage}

Trois artefacts ont été tenus à jour tout au long du projet : le \emph{Product
Backlog}, liste priorisée des fonctionnalités attendues ; le \emph{Sprint Backlog},
sous-ensemble sélectionné pour l'itération en cours ; et l'\emph{incrément}, version
fonctionnelle livrée en fin de sprint. Le backlog complet et le déroulement des six
sprints sont détaillés au chapitre~\ref{chap:analyse}.

Les cérémonies ont été condensées : une réunion de planification en début de sprint,
un point hebdomadaire avec l'encadrement en substitution de la mêlée quotidienne, une
revue de sprint sous forme de démonstration, et une rétrospective écrite. Le suivi
s'est appuyé sur un dépôt Git, dont l'historique des commits documente la progression
réelle et servira à confronter le planning prévu au planning réalisé.

% =============================================================================
\section{Conclusion}
% =============================================================================

Ce chapitre a posé le cadre du projet. Nous avons présenté le Centre National de
l'Informatique et montré comment trois de ses contraintes structurantes — souveraineté
des données, reproductibilité du déploiement, traçabilité — ont orienté les choix
techniques. Nous avons ensuite introduit les notions médicales nécessaires : les cinq
stades du référentiel AASM, l'hypnogramme, le montage polysomnographique et les
classes de sévérité du SAOS.

De ce contexte se dégage une problématique à trois volets : le coût du scorage
manuel, sa variabilité entre experts et l'opacité des modèles qui prétendent
l'automatiser. L'étude de l'existant a montré que les suites commerciales et les
travaux académiques traitent chacun une partie du problème sans le couvrir
entièrement, laissant ouvert l'espace qu'occupe \hypnoria{}.

Le chapitre suivant approfondit le versant scientifique de cette étude : il dresse
l'état de l'art des méthodes de classification automatique des stades du sommeil, des
techniques d'apprentissage ensembliste et des approches d'intelligence artificielle
explicable sur lesquelles s'appuie notre proposition.
